\subsection{实变函数上的复合算子}

设 I 是 R 中包含区间 \(R_{+} = [0, +\infty[\) 的一个区间；设 E 是复数域 C 上的一个向量空间，由定义在 I 上的复值实变函数组成。我们假设对每个 \(a \geqslant 0\) 和每个函数 \(f \in E\) ，函数 \(x \mapsto f(x + a)\) 属于 E；此外，我们假定 E 包含复系数多项式在 I 上的限制以及指数函数 \(e^{\lambda x}\) ，其中 \(\lambda\) 是任意复数。E 到从 I 到复数域 C 的全体映射所成空间中的任何线性映射 U 都称为 E 上的算子；若 \(f \in E\) 且 \(g = U(f)\) ，使用记号

\[
g (x) = U _ {x} ^ {\xi} \left(f (\xi)\right)
\]

其中 \(\xi\) 是右端函数记号中的哑变量（参见 II, p.~58）。对每个 \(a \geqslant 0\) ，把任意函数 \(f \in E\) 对应到函数 \(x \mapsto f(x + a)\) 在 I 上的限制的算子，称为平移 a 的平移算子。

定义 3. 称 E 上的算子 U 为复合算子，如果对每个 \(a \geqslant 0\) ，它与平移 a 的平移算子可交换。

按上面引入的记号，这个定义成为恒等式

\[
U _ {x + a} ^ {\xi} (f (\xi)) = U _ {x} ^ {\xi} (f (\xi + a))\tag{25}
\]

它是关于 \(x\) 和 \(a\) 的恒等式（ \(x \in \mathrm{I}, a \geqslant 0\) ）。可以在该恒等式中交换 \(x\) 与 \(a\) 的角色，只要 \(x \geqslant 0\) ；然后令 \(a = 0\) ；于是对 \(x \geqslant 0\) 得到

\[
U _ {x} ^ {\xi} (f (\xi)) = U _ {0} ^ {\xi} (f (\xi + x))\tag{26}
\]

其中 \(U_{0}\) 是 E 上的线性型，它把每个函数 \(f \in E\) 对应到值 \(g(0)\) ，其中 \(g = U(f)\) 。

若 \(f\) 是多项式，则有 \(f(\xi + x) = \sum_{k=0}^{\infty} \frac{1}{k!} f^{(k)}(\xi)x^k\) ，且公式 (26) 表明 \(U(f)\) 也是多项式；把它限制在 \(x\) 的多项式所成的集合上（系数在 \(\mathbf{C}\) 中，该集合可与代数 \(\mathbf{C}[\mathbf{X}]\) 等同），算子 \(U\) 就是 VI, p.~270 定义 1 意义下的复合算子，从而前面各节的所有结果都可以应用于它。

我们仍把与算子 U 相联系的阿佩尔多项式记作 \(u_{n}\) 。对更一般的函数，与多项式的广义泰勒展开（VI, p.~274, 公式 (13)）相对应，有下列结果：

定理 2. 设 f 是在 I 上具有连续的 \((n+1)^{th}\) 阶导数、并且与其所有导数 \(f^{(m)}\) （ \(1 \leqslant m \leqslant n\) ）一起都属于 E 的函数。若 U 是 E 上一个 \(p \leqslant n\) 阶的复合算子，则对 \(x \geqslant 0\) 和 \(h \geqslant 0\) 有

\[
f ^ {(p)} (x + h) = \sum_ {m = 0} ^ {n} \frac {1}{m !} u _ {m} (x) U _ {h} ^ {\xi} \left(f ^ {(m)} (\xi)\right) + \mathrm{R} _ {n} (x, h)\tag{27}
\]

其中

\[
\mathrm{R} _ {n} (x, h) = - U _ {h} ^ {\xi} \left(\int_ {0} ^ {\xi - x - h} \frac {1}{n !} u _ {n} (x + \eta) f ^ {(n + 1)} (\xi - \eta) d \eta\right)\tag{28}
\]

（广义泰勒展开）。

考虑积分 \(\int_{0}^{\xi-x-h}\frac{1}{n!}u_{n}(x+\eta)f^{(n+1)}(\xi-\eta)d\eta\) （它对所有 \(\xi\in I\) 有定义），并对它应用 \(n+1\) 阶分部积分公式（II, p.~59, 公式 (11)）；利用关系式

\[
u _ {n} ^ {(k)} = n (n - 1) \dots (n - k + 1) u _ {n - k}
\]

（它们是由 (7) （VI, p.~273）通过递推得到的），就得到

\[
\begin{array}{l} \int_ {0} ^ {\xi - x - h} \frac {1}{n !} u _ {n} (x + \eta) f ^ {(n + 1)} (\xi - \eta) d \eta \\ = \sum_ {m = 0} ^ {n} \frac {1}{m !} u _ {m} (x) f ^ {(m)} (\xi) - \sum_ {m = 0} ^ {n} \frac {1}{m !} u _ {m} (\xi - h) f ^ {(m)} (x + h). \end{array}\tag{29}
\]

我们把算子 U 作用于 (29) 的两边（把它们看作 \(\xi\) 的函数），然后取所得函数在变量 \(\xi\) 之值为 h 处的值；注意到，由公式 (26) （VI, p.~277）和 (9) （VI, p.~272），有

\[
U _ {h} ^ {\xi} \big (u _ {m} (\xi - h) \big) = U _ {0} ^ {\xi} \big (u _ {m} (\xi) \big) = \left\{ \begin{array}{l l} 0 & \text {for} m \neq p \\ p! & \text {for} m = p \end{array} \right.
\]

最终就得到 (27)。

